\documentclass[11pt]{article}

\usepackage{amsmath,amssymb}
\usepackage{xurl}
\usepackage[hidelinks]{hyperref}
\setlength{\emergencystretch}{2em}

% Shared commands for notes in docs/paper-gaps/.

\DeclareUnicodeCharacter{2080}{\ifmmode{}_{0}\else\textsubscript{0}\fi}
\DeclareUnicodeCharacter{2081}{\ifmmode{}_{1}\else\textsubscript{1}\fi}
\DeclareUnicodeCharacter{2082}{\ifmmode{}_{2}\else\textsubscript{2}\fi}
\DeclareUnicodeCharacter{2099}{\ifmmode{}_{n}\else\textsubscript{n}\fi}

\newcommand{\Lean}{\textsc{Lean}}
\ExplSyntaxOn
\NewDocumentCommand{\leanid}{m}{
  \ifmmode
    \textsf{\detokenize{#1}}
  \else
    \group_begin:
    \urlstyle{sf}
    \tl_set:Nn \l_tmpa_tl {#1}
    \tl_replace_all:Nnn \l_tmpa_tl {\_} {_}
    \exp_args:NV \path \l_tmpa_tl
    \group_end:
  \fi
}
\ExplSyntaxOff
\providecommand{\tr}{\operatorname{tr}}
\newcommand{\tnleanIssueBase}{https://github.com/LionSR/TNLean/issues/}
\newcommand{\tnleanPrBase}{https://github.com/LionSR/TNLean/pull/}
\newcommand{\tnleanDocBase}{https://sirui-lu.com/TNLean/docs/}
\newcommand{\ghissue}[1]{\href{\tnleanIssueBase#1}{GitHub issue \##1}}
\newcommand{\ghpr}[1]{\href{\tnleanPrBase#1}{PR \##1}}
\newcommand{\doclink}[1]{\href{\tnleanDocBase#1.html}{\path{#1}}}

% Dirac notation and the identity operator, as in blueprint/src/macros/common.tex.
% \providecommand keeps note-local definitions authoritative.
\usepackage{dsfont}
\providecommand{\ket}[1]{|#1\rangle}
\providecommand{\bra}[1]{\langle #1|}
\providecommand{\braket}[2]{\langle #1 | #2 \rangle}
\providecommand{\ketbra}[2]{|#1\rangle\!\langle #2|}
\providecommand{\Id}{\mathds{1}}
\providecommand{\one}{\mathds{1}}
\providecommand{\C}{\mathbb{C}}
\providecommand{\MN}[1]{M_{#1}(\C)}
\providecommand{\MD}{M_D(\C)}

% External referencing for paper sources.  The build helper writes sanitized
% auxiliary files under build/paper-aux/ and excludes duplicated source labels.
\usepackage{xr}

% Import a generated paper auxiliary file with a caller-supplied label prefix.
% The candidate paths support builds from docs/paper-gaps/, the repository root,
% and build/paper-aux/ itself.
\newcommand{\importpaperaux}[2]{%
  \IfFileExists{../../build/paper-aux/#2.aux}{%
    \externaldocument[#1]{../../build/paper-aux/#2}%
  }{%
    \IfFileExists{build/paper-aux/#2.aux}{%
      \externaldocument[#1]{build/paper-aux/#2}%
    }{%
      \IfFileExists{#2.aux}{%
        \externaldocument[#1]{#2}%
      }{}%
    }%
  }%
}

% General external reference macros
\newcommand{\paperref}[2]{\ref{#1-#2}}
\newcommand{\papereqref}[2]{(\ref{#1-#2})}

% CPSV16 source (1606.00608 / MPDO-22-12-17-2.tex) external document and shortcuts
\importpaperaux{cpsv16-}{1606.00608/MPDO-22-12-17-2-tnlean}
\newcommand{\cpsvref}[1]{\paperref{cpsv16}{#1}}
\newcommand{\cpsveqref}[1]{\papereqref{cpsv16}{#1}}

% Verdict marker: \gapnote{<kind>}{<status>}. Kind is the policy
% classification (clarification, local-correction, scope-restriction,
% unfaithful, false-source, open-gap); status is open, wip (elimination
% underway), resolved, or historical. Machine-read by the site index;
% prints a verdict line.
\newcommand{\gapnote}[2]{\par\noindent\textbf{Verdict:} #1 (#2).\par}


\title{Lower-Bound Helper Lemmas for the Theorem~\texttt{thm1} Proof}
\author{}
\date{2026-05-10}

\begin{document}
\maketitle

\section{The source argument}

Cirac--P\'erez-Garc\'ia--Schuch--Verstraete 2016 proves the proportional MPV
Fundamental Theorem (Theorem~\texttt{thm1}) by an argument summarized in
\cite[lines~1170--1192]{Cirac2016MPDO_arXiv}.  The contradiction step at the
heart of the proof is invoked by reference rather than spelled out:

\begin{quote}
    ``Let us first consider $B_k$ for some given $k$.  It is not possible
    that $\langle V^{(N)}(B_k)\,|\,V^{(N)}(A_j)\rangle\to 0$ as
    $N\to\infty$ for all $j$, since otherwise the MPV generated by $A$ and
    $B$ could not be proportional for all $N$ (Lemma~\ref{Lem1}).''
\end{quote}

The cited Corollary~\texttt{Lem1} reads:

\begin{quote}
    ``Any set of NMPV $\{V_j\}_{j=1}^g$ fulfilling
    $\langle V_j\,|\,V_{j'}\rangle\to\delta_{j,j'}$ for $N\to\infty$
    is linearly independent for $N$ sufficiently large.''
    \cite[Corollary~\texttt{Lem1}]{Cirac2016MPDO_arXiv}
\end{quote}

The argument therefore proceeds at a high level: assume the by-contradiction
hypothesis, observe that the joint MPV family is asymptotically orthonormal,
invoke \texttt{Lem1} to conclude eventual linear independence, and read off
that proportionality plus the BNT decompositions force a vanishing
coefficient that the source normalization rules out.

\section{The Lean proof structure}

The Lean development of Theorem~\texttt{thm1} (in
\path{TNLean/MPS/BNT/PermutationRigidity/NonzeroOverlap.lean}) does not
invoke \texttt{Lem1} directly at this step.  In its place it factors the
contradiction through three intermediate lemmas:\footnote{%
\path{TNLean/MPS/BNT/PermutationRigidity/NonzeroOverlap.lean}, lines
near the helpers \leanid{mpvOverlap_total_block_eventually_ge_half},
\leanid{mpvOverlap_total_self_eventually_ge_half}, and
\leanid{mpvOverlap_total_self_le_card}.}

\begin{itemize}
    \item An eventual lower bound
    $\|\langle V^{(N)}(X_{\mathrm{tot}}), V^{(N)}(X_k)\rangle\|\ge 1/2$
    when the coefficient $x_k(N)$ is unit-modulus, the
    other coefficients are bounded by $1$, and the within-family
    overlaps are asymptotically orthonormal.

    \item An eventual lower bound
    $\|\langle V^{(N)}(X_{\mathrm{tot}}), V^{(N)}(X_{\mathrm{tot}})\rangle\|
    \ge 1/2$ when the dominant coefficient is unit-modulus, the others are
    bounded by $1$, and the within-family overlaps are asymptotically
    orthonormal.

    \item A uniform upper bound
    $\|\langle V^{(N)}(X_{\mathrm{tot}}), V^{(N)}(X_{\mathrm{tot}})\rangle\|
    \le M$ from the bounded-coefficient sum.
\end{itemize}

Combined with proportionality and per-$N$ nonzero $c_N$, the lower bounds
on the dominant block's mixed and self overlaps drive the contradiction
without explicit reference to a linear-independence statement.

\section{The mathematical content is the same}

Both the source's \texttt{Lem1}-based route and the Lean lower-bound route
reach the same conclusion through equivalent content.  \texttt{Lem1} is
proved by inverting a Gram matrix that converges to the identity; that
inversion is uniform in $N$ for sufficiently large $N$, which is the same
quantitative statement as the per-$N$ triangle inequality used in the
helpers.  Conversely, given the helper bounds, one recovers eventual
linear independence by the standard determinant lower bound.

\section{Justification for the deviation}

The Lean codebase contains \texttt{Lem1} as
\leanid{eventually_linearIndependent_of_overlap_tendsto_orthonormal}
\cite[\path{TNLean/MPS/BNT/Basic.lean:91}]{}.\footnote{%
The Lean form takes a finite family $A : \mathrm{Fin}\ g \to
\mathrm{MPSTensor}\ d\ (\mathrm{dim}\ j)$ together with the asymptotic
self-overlap and off-diagonal-overlap hypotheses, and concludes
$\forall^{\infty}\,N$, $\{V^{(N)}(A_j)\}_j$ is linearly independent in
$\mathrm{MPVSpace}\ d\ N$.}  Using it directly at the contradiction step
would route the proof through linear-independence reasoning over a joint
family with an additional element $V^{(N)}(B_k)$.  Translating
``coefficient extraction from a linear relation in an asymptotically LI
family'' into per-$N$ statements is more involved in Lean than the direct
triangle-inequality bounds carried by the helpers.  The helpers also
expose the key quantitative inputs (the dominant-block normalization and
the uniform sub-dominant bounds) at the surface where the contradiction
becomes visible, which keeps the rest of the proof readable.

The deviation is therefore a proof-style choice, not a strengthening of
hypotheses or a weakening of conclusions.  The helpers add no constraint
on the canonical-form data beyond what \texttt{Lem1} already requires,
and the contradiction they support is the same one
\cite{Cirac2016MPDO_arXiv} draws.

\section{Conclusion}

The three lower-bound helper lemmas in the Lean proof of
Theorem~\texttt{thm1} are not stated in
\cite{Cirac2016MPDO_arXiv}; the paper instead invokes
Corollary~\texttt{Lem1} abstractly at the corresponding step.  The
mathematical content of the helpers is a concrete per-$N$ derivative of
\texttt{Lem1}'s consequences, and the eventual existence of an
asymptotically orthonormal coefficient family is exactly what
\texttt{Lem1} would provide.  The Lean route is preferred for
formalization tractability; the deviation is structural rather than
mathematical.

\bibliographystyle{alpha}
\bibliography{references}

\end{document}
